%% ma: L12-B19-0035
%% lop: 12
%% chuong: 6
%% bai: 19
%% dang: 12.19.1
%% loai: TL
%% muc: TH
%% dap_an: "-"
%% nguon: "(NBV)-10 ĐỀ KIỂM TRA HK2-TOÁN 12-DUNG CHUNG 3 SÁCH.pdf (H03, Đề số 3), Phần 4 Câu 3 (tự luận)"
%% kien_thuc: "Bài 19 (công thức xác suất toàn phần)"
%% trang_thai: chinh-thuc
%% ghi_chu: "Đáp án gốc 0,625 khớp. Đề gốc có thể hiểu mập mờ cụm ‘báo sai người vô tội’ (đã hiểu là P(B | Ā) = 0,1 theo lời giải gốc)."
\begin{ex}
Một cơ quan điều tra, người ta dùng máy dò tìm tội phạm, kinh nghiệm cho biết cứ $10$ người bị tình nghi thì có $7$ người là tội phạm. Máy báo đúng người có tội với xác suất là $0{,}85$. Máy báo sai người vô tội với xác suất là $0{,}1$. Một người được máy phân tích, hãy tính xác suất máy báo người này là tội phạm.
\loigiai{Gọi $A$: ``người bị tình nghi là tội phạm'', $B$: ``máy báo người này là tội phạm''. Khi đó $P(A)=0{,}7$, $P\left(\overline A\right)=0{,}3$, $P(B\mid A)=0{,}85$, $P\left(B\mid\overline A\right)=0{,}1$ (máy báo sai người vô tội, tức báo người vô tội là tội phạm).

Theo công thức xác suất toàn phần: $P(B)=P(A)P(B\mid A)+P\left(\overline A\right)P\left(B\mid\overline A\right)=0{,}7\cdot0{,}85+0{,}3\cdot0{,}1=0{,}625$.}
\end{ex}
